% Pakete und Einstellungen – wird von main.tex eingebunden
\usepackage[T1]{fontenc}
\usepackage{lmodern}
\usepackage[ngerman]{babel}
\usepackage{csquotes}
\usepackage[a4paper,margin=2.5cm]{geometry}
\usepackage{microtype}
\usepackage{amsmath,amssymb}
\usepackage{siunitx}
\usepackage{booktabs}
\usepackage{graphicx}
\usepackage{caption}
\usepackage{subcaption}
\usepackage{tikz}
\usetikzlibrary{arrows.meta,calc,positioning}
% Farben der Zeigerdiagramme
\colorlet{cU}{blue!75!black}\colorlet{cI}{red!80!black}\colorlet{cD}{gray}
\colorlet{cP}{green!50!black}\colorlet{cQ}{orange!90!black}
\usepackage{pgfplots}
\pgfplotsset{compat=1.16}
% Offene Punkte als sichtbare Boxen. Vor der Abgabe:  [disable]  statt [textsize=footnotesize]
\usepackage[disable]{todonotes}
\usepackage[hidelinks]{hyperref}

\graphicspath{{img/}}
\captionsetup{font=small,labelfont=bf}

% Dezimalpunkt (Schweizer Schreibweise), Einheiten
\sisetup{output-decimal-marker={.},per-mode=symbol}
\DeclareSIUnit{\VAr}{var}

% Abkürzungen für Formelzeichen
\newcommand{\Udc}{U_{\mathrm{dc}}}
\newcommand{\UN}{U_{\mathrm{N}}}
\newcommand{\UWR}{U_{\mathrm{WR}}}
\newcommand{\uN}{u_{\mathrm{N}}}
\newcommand{\uWR}{u_{\mathrm{WR}}}
\newcommand{\iN}{i_{\mathrm{N}}}
\newcommand{\Th}{\vartheta}

% Gemeinsamer Stil für Diagramme
\pgfplotsset{
  rep/.style={width=0.92\linewidth,height=5.2cm,grid=major,
    tick label style={font=\small},label style={font=\small},
    legend style={font=\small},unbounded coords=jump}
}
